\documentclass[11pt,a4paper]{article}
\usepackage[top=2.0cm,bottom=2.0cm,left=2.2cm,right=2.2cm]{geometry}
\usepackage[T1]{fontenc}
\usepackage[utf8]{inputenc}
\usepackage{xcolor}
\usepackage{hyperref}

\pagestyle{empty}
\setlength{\parskip}{0.65em}
\setlength{\parindent}{0pt}

\definecolor{accent}{HTML}{1F4E78}
\hypersetup{colorlinks=true, urlcolor=accent, linkcolor=accent}

\begin{document}

\textbf{\large Prabhavit Gupta}\\
Via Zanetti 10, 30175 Venice, Italy ~$\vert$~ +39 351 610 7807 ~$\vert$~ \href{mailto:prabhavitg@gmail.com}{prabhavitg@gmail.com}\\
\href{https://www.linkedin.com/in/prabhavit/}{linkedin.com/in/prabhavit} ~$\vert$~ \href{https://github.com/prab-gupta}{github.com/prab-gupta}\\
26 September 2026

\textbf{Hiring Team}\\
MC Engineering\\
Application for Gen AI Engineer -- Italy (Remote)

Dear Hiring Manager,

I am writing to express my strong interest in the Gen AI Engineer position at MC Engineering. As an engineer specializing in Generative AI, LLM orchestration, and intelligent automation based in Venice, I have followed MC Engineering's consulting engagements delivering digital transformation and advanced technology integration. Your focus on integrating generative models, autonomous AI agents, and intelligent workflows into enterprise microservices and platforms aligns directly with my hands-on background in LLM pipelines, rapid prototyping, and complex API integrations.

Over the past three years, I have successfully transformed abstract business requirements into production-ready features across multiple European technology projects spanning different countries and sectors; from a sports management platform in Luxembourg to strategic consultant platforms in France and Spain. Operating extensively as a freelance AI Solutions Developer, I engineered AI-based data parsing tools using n8n for UltraStrategy, automated complex workflows for Get2Germany, and executed rigorous quality assurance and API integrations for Clubee.

During my tenure at FlowSpace (London, remote), I led the migration of a live marketplace from a legacy stack to a modern Next.js and Laravel architecture, scaling the platform from 2,000 to over 4,500 active users. On the AI side, I architected an AI document ingestion pipeline using TypeScript, Claude, and n8n that parses unstructured brochures into structured listings with human-in-the-loop review. I also engineered customer support AI agents connected to platform data with confidence-based escalation, and integrated third-party booking APIs (OfficeRnD, Nexudus) managing automated synchronization.

My daily engineering workflow centers on practical enterprise GenAI integration: designing prompt chaining sequences, RAG architectures, and workflow pipelines (Python, TypeScript, n8n) that connect external REST APIs and enterprise backends with foundational models like Anthropic's Claude and OpenAI. Leveraging AI coding tools like Claude Code alongside modern engineering practices, I prioritize deterministic validation, robust error handling, and measurable business impact over brittle proof-of-concepts.

As a graduate of Ca' Foscari University of Venice, I hold an Italian European Long-Term Residence Permit (EU citizen equivalent for work authorization, requiring no visa sponsorship) and communicate fluently in English (C2). Living in Italy, I am comfortable collaborating with Italian-speaking teams and fully available for remote work with travel to client sites or MC Engineering offices in Turin, Milan, or Rome as project demands.

I would welcome the opportunity to discuss how my practical experience in GenAI engineering, LLM orchestration, and enterprise integration can contribute to MC Engineering's client engagements.

Thank you for your time and consideration.

Sincerely,\\
Prabhavit Gupta

\vspace{0.8em}
{\footnotesize \textit{I hereby authorize the processing of my personal data contained in this document for the purposes of recruitment and selection in compliance with the General Data Protection Regulation (GDPR) and Italian Legislative Decree 196/2003, as modified by Legislative Decree 101/2018.}\par}

\end{document}
